\chapter{Causal discovery}
\label{cha:1}

\section{Introduction to Causality}

Understanding causal relationships between variables is a fundamental objective in many scientific and engineering disciplines. While statistical analysis often reveals correlations between variables, correlation alone does not imply causation. Two variables may appear related because one directly influences the other, because they share a common cause, or because their relationship arises through indirect interactions within a larger system. Distinguishing between these possibilities is essential when the goal is to understand how changes in one variable influence another \cite{Pearl2009,Spirtes2000}.

Causal reasoning is particularly important when decisions or interventions must be made based on data. In many real-world situations, experiments that directly manipulate variables are difficult, expensive, or ethically constrained. As a result, researchers and engineers often rely on observational data to infer how systems behave. Causal analysis aims to extract cause--effect relationships from such data in order to support prediction, explanation, and decision making \cite{Pearl2009}.

One common way to represent causal relationships is through graphical models, where variables are represented as nodes and causal influences are represented as directed edges. These graphical representations allow complex systems to be described in a structured way and form the basis of many causal discovery methods. By combining graphical representations with statistical information about dependencies between variables, it becomes possible to infer aspects of the underlying causal structure from observed data \cite{Spirtes2000}.

The study of causal inference and causal discovery has developed rapidly over the past decades and is now an active research field spanning statistics, computer science, economics, and the natural sciences. Several theoretical frameworks have been developed to formalize causal reasoning and to provide algorithms that can infer causal relations from data under certain assumptions \cite{Pearl2009,Spirtes2000,Peters2017}.
